\documentclass[journal]{IEEEtran}

\usepackage{amsmath,amssymb,amsfonts,mathtools}
\usepackage{graphicx}
\usepackage{booktabs}
\usepackage{tabularx}
\usepackage{array}
\usepackage{cite}
\usepackage{xcolor}
\usepackage[colorlinks=true,linkcolor=blue!60!black,citecolor=blue!60!black,urlcolor=blue!60!black]{hyperref}

\newcolumntype{L}[1]{>{\raggedright\arraybackslash}p{#1}}
\newcolumntype{C}[1]{>{\centering\arraybackslash}p{#1}}
\newcolumntype{Y}{>{\raggedright\arraybackslash}X}

\begin{document}

\title{Detecting Misinterpreted but Genuine Citations Through Automated Claim Verification}

\author{\IEEEauthorblockN{Tushar Kumar}
\IEEEauthorblockA{Department of Information Technology\\
Indian Institute of Information Technology, Allahabad\\
Prayagraj, Uttar Pradesh, India\\
Email: mhm2026010@iiita.ac.in}}

\maketitle

\begin{abstract}
The accelerating volume of scholarly publishing and the widespread adoption of generative writing assistants have amplified a subtle yet pervasive threat to scientific integrity: \emph{misinterpreted but genuine citations}. Unlike fabricated or phantom references, a misinterpreted genuine citation points to an authentic, peer-reviewed publication with valid Digital Object Identifier (DOI) metadata and topical relevance, yet the citing assertion exaggerates, overgeneralizes, conflates correlation with causation, distorts numerical metrics, or directly contradicts the cited source's actual findings. Conventional bibliographic validators verify only catalog existence, whereas existing citation-intent classifiers rely solely on the citing author's local phrasing and current scientific fact-checking benchmarks restrict evaluation to isolated sentences matched against paper abstracts. This paper presents CEDAR (Claim-Evidence Decomposition, Attribution, and Risk-Calibrated Verification), a methodological framework for detecting misinterpreted genuine citations across full-text scholarly manuscripts. First, we formulate a cohesion-penalized discourse window optimization and a bipartite claim-to-anchor projection that decomposes multi-clause citing sentences into decontextualized atomic propositions paired with exact bibliographic targets. Second, we design a section-weighted multimodal Reciprocal Rank Fusion (RRF) and cross-encoder reranking architecture that retrieves evidentiary rationales across narrative prose, structured tables, and figure captions. Third, we move beyond rigid tri-state natural language inference (NLI) by introducing a continuous Epistemic Modality Delta ($\Delta_{\mathcal{E}}$), scope-divergence scoring, and neuro-symbolic numerical tolerance verification. Finally, to prevent uncalibrated false accusations in editorial workflows, we formulate uncertainty-aware decision routing via split conformal prediction sets with finite-sample marginal coverage guarantees. We conclude with a reproducible, leakage-free experimental evaluation protocol and an analysis of socio-technical deployment barriers.
\end{abstract}

\begin{IEEEkeywords}
Citation Verification, Claim Verification, Scientific Natural Language Processing, Natural Language Inference, Multimodal Evidence Retrieval, Epistemic Calibration, Conformal Prediction, Citation Integrity.
\end{IEEEkeywords}

\section{Introduction}
\IEEEPARstart{S}{cholarly} communication relies on the normative assumption that an inline citation provides verifiable evidentiary grounding for the proposition to which it is attached. Across empirical and theoretical disciplines, citations trace intellectual lineage, substantiate methodological choices, bound experimental assumptions, and anchor new claims within peer-reviewed findings \cite{moravcsik1975some,jurgens2018measuring}. In recent years, however, two concurrent pressures have strained the reliability of this mechanism: the rapid growth of global publication volume, which limits the time peer reviewers can devote to auditing primary sources, and the proliferation of large language models (LLMs) in academic drafting workflows \cite{hosseini2024using,thorp2023chatgpt}. While early discussions of generative AI in science focused primarily on outright reference fabrication---commonly termed ``hallucinated'' or phantom citations---empirical bibliometric audits demonstrate that an equally damaging failure mode is pervasive even in human-authored literature: \emph{misinterpreted but genuine citations} \cite{delacey1985accurate,greenberg2009citation,jergas2015quotation,mogull2017accuracy}.

\begin{table*}[t]
\centering
\footnotesize
\caption{Stratified Epistemic Dimensions of Scholarly Citation Integrity and Verification}
\label{tab:integrity_dimensions}
\renewcommand{\arraystretch}{1.15}
\begin{tabularx}{\textwidth}{L{2.5cm} L{4.2cm} L{4.2cm} Y}
\toprule
\textbf{Verification Layer} & \textbf{Core Epistemic Question} & \textbf{Primary Failure Mode} & \textbf{Representative Computational Audit Mechanism} \\
\midrule
1. Source Authenticity & Does the cited work genuinely exist as an indexed scholarly record? & Fabricated references, LLM-hallucinated phantom papers \cite{hosseini2024using,thorp2023chatgpt}. & Crossref, PubMed, and OpenAlex DOI resolution; digital catalog verification. \\
2. Bibliographic Fidelity & Do the reference string fields match the official publication record? & Mutated author lists, wrong volume or pagination, corrupted DOIs \cite{simkin2003read,harzing2002referencing}. & Structured header/reference parsing via GROBID \cite{lopez2009grobid} and fuzzy record linkage. \\
3. Topical Relevance & Does the cited document address the subject domain of the citing passage? & Perfunctory or honorary citation padding; citation metric gaming \cite{moravcsik1975some}. & Sparse lexical matching (BM25 \cite{robertson2009probabilistic}) and citation-graph document embeddings (SPECTER \cite{cohan2020specter}). \\
4. Propositional Entailment & Does the full text of the cited work substantiate the atomic claim? & Unsubstantiated leaps, missing evidence, or directional contradiction \cite{wadden2020fact,wadden2022multivers}. & Discourse-scoped claim decomposition \cite{kamoi2023wice,min2023factscore}, hybrid retrieval, and cross-encoder NLI. \\
5. Epistemic \& Scope Fidelity & Are hedging qualifiers, experimental bounds, and quantitative metrics preserved? & Speculative elevation, causal inflation, animal-to-human scope jump, numerical distortion \cite{greenberg2009citation,hyland1998hedging,mogull2017accuracy}. & Epistemic modality delta ($\Delta_{\mathcal{E}}$), table/chart verification \cite{lu2023scitab,akhtar2023reading}, and conformal risk calibration \cite{angelopoulos2023conformal,sadinle2019least}. \\
\bottomrule
\end{tabularx}
\end{table*}

A misinterpreted but genuine citation occurs when a citing manuscript references an authentic, cataloged, and topically relevant publication possessing valid bibliographic metadata, yet the in-text assertion attached to that citation is unsupported, exaggerated, overgeneralized, stripped of essential boundary conditions, or directly contradicted by the primary source \cite{greenberg2009citation,rekdal2014sinking,pavlovic2021accurate}. Because the cited work genuinely exists and its Digital Object Identifier (DOI) resolves cleanly in bibliographic databases such as Crossref, PubMed, or OpenAlex, automated editorial screening tools classify the reference as valid. In effect, the prestige and bibliographic authenticity of the cited publication act as a veneer of legitimacy that masks an unsubstantiated or distorted scientific claim.

\subsection{Authenticity Versus Epistemic Support}
To formalize this vulnerability, it is necessary to separate scholarly citation integrity into five distinct epistemic layers, summarized in Table~\ref{tab:integrity_dimensions}. Automated reference managers, publisher XML validators, and bibliographic parsers such as GROBID \cite{lopez2009grobid} operate exclusively at Layers 1 and 2 (Source Authenticity and Bibliographic Fidelity). Citation recommendation engines and document similarity models such as SPECTER \cite{cohan2020specter} operate at Layer 3 (Topical Relevance). However, a reference can satisfy all three lower layers---existing in the catalog, having exact metadata, and sharing high vocabulary overlap with the citing paragraph---while failing at Layers 4 and 5 (Propositional Entailment and Epistemic/Scope Fidelity).

Empirical meta-analyses across medicine, biology, ecology, social sciences, and general science confirm that failures at Layers 4 and 5 are widespread. In a foundational audit of six medical journals, de Lacey et al.\ \cite{delacey1985accurate} distinguished between bibliographic citation errors (24\% overall, 8\% major) and \emph{quotation errors}---where the citing text misrepresented or failed to reflect the original author's findings---reporting a 15\% quotation error rate. A comprehensive systematic review and meta-analysis of 28 studies covering 7,321 medical citations by Jergas and Baethge \cite{jergas2015quotation} estimated a pooled major quotation error rate of 11.9\% (with total quotation errors reaching 24.4\%), while Mogull \cite{mogull2017accuracy} recalculated medical quotation error rates at 14.5\% after controlling for indirect citations. Similar patterns have been documented in marine biology, where Todd et al.\ \cite{todd2010marine} found that 24.2\% of citations failed to unambiguously support the cited assertion, in frequently cited biomedical articles audited by Pavlovic et al.\ \cite{pavlovic2021accurate}, and in general science and humanities journals examined by Smith and Cumberledge \cite{smith2020quotation}.

\subsection{Citation Contagion and Cascading Scientific Risk}
Once a misinterpreted citation enters the peer-reviewed literature, it rarely remains isolated. Using a stochastic model of duplicated typographic errors in reference lists, Simkin and Roychowdhury \cite{simkin2003read} estimated that approximately 78\% of citing authors copy references from secondary sources without reading the original publication. Harzing \cite{harzing2002referencing} and Rekdal \cite{rekdal2014sinking,rekdal2014urban} traced how secondhand citation copying transforms tentative hypotheses or misquoted statistics into entrenched ``academic urban legends.'' In a landmark network analysis of 242 publications and 675 citations investigating the hypothesis that $\beta$-amyloid protein is pathologically produced inside skeletal muscle fibers in inclusion body myositis, Greenberg \cite{greenberg2009citation} demonstrated how speculative discussion remarks were repeatedly cited as definitive empirical proof, while studies refuting the claim were systematically ignored---a phenomenon Greenberg termed \emph{unfounded authority} and \emph{citation distortion}. In preclinical and clinical research, where reproducibility is already fragile \cite{ioannidis2005why,begley2012raise}, such citation cascades can misdirect experimental funding, distort clinical guidelines, and propagate erroneous benchmark baselines.

\subsection{Contributions of This Work}
Despite rapid progress in automated fact-checking \cite{thorne2018fever,zeng2021automated} and scientific natural language processing (NLP) \cite{wadden2020fact,vladika2023scientific}, existing verification pipelines remain ill-equipped to audit misinterpreted genuine citations in full-text manuscripts. To bridge this gap, this paper develops a rigorous methodological framework---CEDAR (\textbf{C}laim-\textbf{E}vidence \textbf{D}ecomposition, \textbf{A}ttribution, and \textbf{R}isk-Calibrated Verification)---making four primary contributions:
\begin{enumerate}
\item \textbf{Problem Formalization and Grounded Taxonomy:} We formalize the distinction between bibliographic validity, topical similarity, and discourse-scoped epistemic entailment, and establish an eight-class taxonomy of citation misinterpretation grounded in verified empirical audits \cite{delacey1985accurate,greenberg2009citation,jergas2015quotation,mogull2017accuracy}.
\item \textbf{Discourse-Scoped Bipartite Claim Decomposition:} We formulate an optimization objective for extracting adaptive citation discourse windows and projecting multi-clause hypotactic citing sentences into decontextualized atomic claim--anchor pairs $(c_k, a_i)$, eliminating cross-reference conflation in compound sentences \cite{kamoi2023wice,min2023factscore}.
\item \textbf{Multimodal Section-Weighted Evidence Retrieval and Epistemic Delta Verification:} We integrate lexical retrieval (BM25 \cite{robertson2009probabilistic}), dense citation-informed embeddings (SPECTER \cite{cohan2020specter}), late-interaction token matching (ColBERT \cite{khattab2020colbert}), and structured table/chart parsing (TaPas \cite{herzig2020tapas}, SCITAB \cite{lu2023scitab}, chart reasoning \cite{akhtar2023reading}) via section-weighted Reciprocal Rank Fusion \cite{cormack2009reciprocal}. We further introduce a continuous \emph{Epistemic Modality Delta} ($\Delta_{\mathcal{E}}$) and scope divergence metric to detect hedging erasure and causal inflation \cite{hyland1998hedging}.
\item \textbf{Conformal Risk Control and Reproducible Evaluation Protocol:} We formulate distribution-free split conformal prediction sets \cite{sadinle2019least,angelopoulos2023conformal} that guarantee finite-sample marginal coverage $1-\alpha$, routing ambiguous cases to human auditors via explainable Verification Cards, and specify a complete, leakage-free experimental evaluation protocol.
\end{enumerate}

The remainder of this paper is organized as follows. Section~\ref{sec:related} critically reviews related work across scientometrics, citation analysis, and scientific fact verification. Section~\ref{sec:gap_problem} defines the research gap, mathematical problem formulation, research questions, and objectives. Section~\ref{sec:taxonomy} presents the taxonomy of citation misinterpretation, distinguishing verified historical cases from illustrative diagnostic scenarios. Section~\ref{sec:framework} details the 12-stage CEDAR architecture. Section~\ref{sec:math} establishes the mathematical formulations and algorithmic design. Section~\ref{sec:llm_critique} analyzes the failure modes of LLMs in citation auditing. Section~\ref{sec:evaluation} details the proposed reproducible evaluation methodology. Section~\ref{sec:limitations} discusses limitations, threats to validity, and future directions, and Section~\ref{sec:conclusion} concludes the paper.

\section{Background and Critical Literature Review}
\label{sec:related}

\subsection{Scientometrics of Quotation Errors and Citation Cascades}
Early scientometric research by Moravcsik and Murugesan \cite{moravcsik1975some} demonstrated that nearly 40\% of citations in theoretical high-energy physics were perfunctory, serving to acknowledge broad background rather than supply necessary methodological or empirical premises. Subsequent audits separated citation errors into two orthogonal categories: \emph{bibliographic errors} (garbled metadata) and \emph{quotation errors} (content misrepresentation) \cite{delacey1985accurate,jergas2015quotation}. Across clinical medicine \cite{delacey1985accurate,mogull2017accuracy,pavlovic2021accurate}, ecology and marine biology \cite{todd2010marine}, management science \cite{harzing2002referencing}, and general science \cite{smith2020quotation}, empirical audits consistently show that 12\% to 25\% of citations fail to substantiate the specific proposition for which they are cited. Furthermore, Rekdal \cite{rekdal2014sinking,rekdal2014urban} and Greenberg \cite{greenberg2009citation} documented that quotation errors are predominantly directional: caveats, negative controls, and linguistic hedges \cite{hyland1998hedging} are stripped away as assertions propagate through secondary and tertiary citations.

\subsection{Citation Context Analysis and Intent Classification}
In computational linguistics, Teufel et al.\ \cite{teufel2006automatic} pioneered automated classification of citation function, identifying whether a citing sentence contrasts results, adopts a method, or establishes background. Jurgens et al.\ \cite{jurgens2018measuring} expanded this paradigm into six functional citation frames to model disciplinary evolution. At production scale, Nicholson et al.\ \cite{nicholson2021scite} introduced \emph{scite}, training deep learning classifiers on millions of citation contexts (``citances'') to label citations as \emph{supporting}, \emph{mentioning}, or \emph{contrasting}.

Crucially, however, citation intent classification is an \emph{intra-document} rhetorical task rather than an \emph{inter-document} epistemic verification task. Platforms such as scite \cite{nicholson2021scite} classify the rhetorical stance expressed by the \emph{citing} author within the citing snippet. If a citing author writes, ``Prior clinical trials by Author X conclusively demonstrated a 50\% mortality reduction [j],'' a citation-intent classifier labels the citation as \emph{supporting} based on the phrase ``conclusively demonstrated,'' even if Reference~$[j]$ actually reported a non-significant trend or an animal model study. Detecting misinterpretation requires conditioning simultaneously on the citing context and the full evidentiary content of the cited source.

\begin{table*}[t]
\centering
\footnotesize
\caption{Critical Comparison of Existing Fact-Checking, Citation Analysis, and Scientific Verification Paradigms with CEDAR}
\label{tab:literature_comparison}
\renewcommand{\arraystretch}{1.14}
\begin{tabularx}{\textwidth}{L{2.6cm} L{2.2cm} C{1.4cm} C{1.5cm} C{1.5cm} C{1.5cm} C{1.4cm} Y}
\toprule
\textbf{System / Benchmark} & \textbf{Primary Domain \& Source Granularity} & \textbf{Discourse Scoping} & \textbf{Bipartite Claim Split} & \textbf{Full-Text \& Tables} & \textbf{Epistemic Delta ($\Delta_{\mathcal{E}}$)} & \textbf{Conformal Bound} & \textbf{Key Limitation for Auditing Misinterpreted Genuine Citations} \\
\midrule
FEVER \cite{thorne2018fever} & Wikipedia articles (introductory text) & No (single sentence) & No & No (text only) & No & No & Synthetic crowd-written claims lack scientific hedging, multi-ref syntax, and quantitative metrics. \\
scite \cite{nicholson2021scite} & Scholarly citances (citing text only) & Local citance window & No & No (ignores source text) & No & No & Classifies rhetorical intent of the citing author; does not verify entailment against the cited paper. \\
SciFact \cite{wadden2020fact} & Biomedical claims vs.\ PubMed abstracts & No (rewritten claims) & No & No (abstracts only) & No & No & Abstracts omit ablation tables, negative sub-results, and specific baseline numbers cited in full text. \\
COVID-Fact \cite{saakyan2021covidfact} & COVID-19 claims vs.\ lit./media snippets & No & No & Partial text snippets & No & No & Focuses on binary supported/refuted claims without modeling scope extrapolation or modality shift. \\
MultiVerS \cite{wadden2022multivers} & Biomedical claims vs.\ full-document text & No & No & Full text (prose only) & No & No & Encodes long prose paragraphs via Longformer but ignores structured tables, charts, and multi-anchor clauses. \\
WiCE \cite{kamoi2023wice} / FActScore \cite{min2023factscore} & Wikipedia / LLM long-form biographies & Single sentence & Sub-claim split & Web / Wiki prose & No & No & Decomposes propositions but lacks bipartite citation-anchor binding and scientific modality/scope calibration. \\
ALCE \cite{gao2023enabling} / Rashkin \cite{rashkin2023measuring} & LLM-generated QA with citations & Prompt context & Sentence-level NLI & Retrieved web passages & No & No & Evaluates generative attribution via off-the-shelf NLI without numerical or conformal risk control. \\
SCITAB \cite{lu2023scitab} / Akhtar \cite{akhtar2023reading} & Scientific tables and chart images & No & No & Isolated tables/charts & No & No & Operates on pre-isolated tables or charts rather than end-to-end full-text citation-to-source verification. \\
\textbf{CEDAR (Proposed)} & \textbf{Full-text citing \& cited manuscripts} & \textbf{Cohesion window} & \textbf{Bipartite $(c_k, a_i)$} & \textbf{Prose + Tables + Captions} & \textbf{Yes ($\Delta_{\mathcal{E}}, D_{\text{scope}}$)} & \textbf{Yes ($1-\alpha$ bound)} & \textbf{Requires full-text access ($h_{\text{doc}}=1$) and calibrated domain-adapted NLI backbones.} \\
\bottomrule
\end{tabularx}
\end{table*}

\subsection{Automated Fact Verification and Scientific NLI}
The modern paradigm of automated fact verification was standardized by the RTE challenges \cite{dagan2006pascal}, broad-coverage NLI corpora such as MultiNLI \cite{williams2018broad}, and the FEVER benchmark \cite{thorne2018fever,zeng2021automated}, which established a three-stage pipeline: document retrieval, sentence selection, and three-way entailment classification (\textsc{Supported}, \textsc{Refuted}, \textsc{NotEnoughInfo}). However, as demonstrated by Wadden et al.\ in SciFact \cite{wadden2020fact} and surveyed by Vladika and Matthes \cite{vladika2023scientific}, general-domain NLI models degrade substantially on scientific claims due to specialized terminology, complex subordination, and quantitative comparisons. Domain-specific corpora including PubMedQA \cite{jin2019pubmedqa}, COVID-Fact \cite{saakyan2021covidfact}, and zero-shot claim generation frameworks \cite{wright2022generating} advanced biomedical claim verification, while MultiVerS \cite{wadden2022multivers} incorporated weak supervision over full-document prose.

More recently, research on attribution in natural language generation---formalized under the Attributable to Identified Sources (AIS) framework by Rashkin et al.\ \cite{rashkin2023measuring} and operationalized for LLM citation generation by Gao et al.\ (ALCE) \cite{gao2023enabling}---and fine-grained claim decomposition via WiCE \cite{kamoi2023wice} and FActScore \cite{min2023factscore} have demonstrated that decomposing complex sentences into atomic propositions substantially improves verification fidelity. Simultaneously, specialized benchmarks for scientific tables (TaPas \cite{herzig2020tapas}, SCITAB \cite{lu2023scitab}) and chart images \cite{akhtar2023reading} have shown that crucial empirical evidence frequently resides outside narrative prose.

\section{Research Gap, Problem Formulation, and Research Questions}
\label{sec:gap_problem}

\subsection{Synthesis of the Research Gap}
A critical comparison of the literature (Table~\ref{tab:literature_comparison}) reveals four specific, verifiable structural gaps in existing scientific verification systems:
\begin{enumerate}
\item \textbf{Multi-Anchor Syntactic Conflation and Discourse Truncation:} Existing scientific verification benchmarks (SciFact \cite{wadden2020fact}, MultiVerS \cite{wadden2022multivers}) evaluate pre-extracted, standalone claims against a single target document. In authentic manuscripts, however, citing sentences frequently contain multiple clauses anchored to different references, or depend on anaphoric pronouns whose antecedents lie in preceding sentences \cite{jurgens2018measuring,kamoi2023wice}. Evaluating a raw citing sentence against a single reference conflates claims meant for other citations, inducing severe false-positive refutations.
\item \textbf{The Abstract-Only and Prose-Only Bottleneck:} Most scientific fact-checking systems verify claims solely against paper abstracts \cite{wadden2020fact,jin2019pubmedqa}. Yet authors routinely cite papers from corpora such as S2ORC \cite{lo2020s2orc} for specific experimental baselines, ablation studies, or hyperparameter settings reported only in full-text sections, structured data tables \cite{herzig2020tapas,lu2023scitab}, or figure plots \cite{akhtar2023reading}.
\item \textbf{Insensitivity to Epistemic Modality and Scope Drift:} Standard NLI formulations \cite{williams2018broad,thorne2018fever} treat entailment as a discrete categorical decision. They lack explicit mathematical mechanisms to penalize \emph{epistemic overclaiming}---where a hedged hypothesis (``may suggest a potential link'' \cite{hyland1998hedging}) is cited as an established causal fact \cite{greenberg2009citation}---or \emph{scope extrapolation}, where findings from a restricted cohort (e.g., murine \emph{in vitro} cells) are generalized to human clinical efficacy \cite{begley2012raise,mogull2017accuracy}.
\item \textbf{Uncalibrated Confidence and High-Stakes False Accusation Risk:} Neural NLI classifiers and generative LLMs exhibit severe miscalibration and sycophancy on out-of-distribution scientific text \cite{liu2024lost,angelopoulos2023conformal}. Existing citation checkers lack distribution-free statistical coverage guarantees to bound false accusations before flagging a manuscript for editorial review \cite{sadinle2019least}.
\end{enumerate}

\subsection{Formal Problem Statement}
Let $D_{\text{cite}}$ denote a citing scholarly manuscript composed of a sequence of sentences $\mathcal{S} = (s_1, s_2, \dots, s_M)$ and a structured bibliography $\mathcal{B} = \{b_1, b_2, \dots, b_J\}$. Within $D_{\text{cite}}$, let $\mathcal{A} = \{a_1, a_2, \dots, a_I\}$ denote the set of inline citation anchors occurring in the text, where each anchor $a_i \in \mathcal{A}$ appears inside a focal sentence $s_{t(i)} \in \mathcal{S}$ and maps via a deterministic bibliographic resolution function $\mu: \mathcal{A} \to \mathcal{B}$ to a bibliography entry $b_j = \mu(a_i)$.

When $b_j$ resolves to a genuine, cataloged publication $D_{\text{src}}^{(j)}$ with verified metadata (i.e., Layer 1 and Layer 2 checks pass), we represent $D_{\text{src}}^{(j)}$ as a multimodal collection of evidentiary units:
\begin{equation}
D_{\text{src}}^{(j)} = \mathcal{U}_{\text{prose}}^{(j)} \cup \mathcal{U}_{\text{table}}^{(j)} \cup \mathcal{U}_{\text{fig}}^{(j)} = \{e_1, e_2, \dots, e_{N_j}\},
\label{eq:source_units}
\end{equation}
where each unit $e_r \in D_{\text{src}}^{(j)}$ belongs to a rhetorical section $z(e_r) \in \mathcal{Z}$ ($\text{Title}$, $\text{Abstract}$, $\text{Intro}$, $\text{Methods}$, $\text{Results}$, $\text{Discussion}$, $\text{Table}$, or $\text{Caption}$). Given a citation anchor $a_i$ and its resolved source $D_{\text{src}}^{(j)}$, the automated citation verification problem is to: (i) extract the minimal discourse context $C(a_i)$ surrounding $s_{t(i)}$; (ii) decompose $C(a_i)$ into decontextualized atomic propositions $\mathcal{C}(a_i) = \{c_{i,1}, \dots, c_{i,m_i}\}$ specifically attributed to $a_i$; (iii) retrieve a minimal sufficient evidentiary rationale set $E^*(c_{i,k}) \subseteq D_{\text{src}}^{(j)}$; and (iv) compute a calibrated verification tuple:
\begin{equation}
\mathcal{V}(c_{i,k}, D_{\text{src}}^{(j)}) = \Big(\hat{y}_{i,k},\; \mathcal{C}_{\alpha}(c_{i,k}, E^*),\; \Delta_{\mathcal{E}}(c_{i,k}, E^*),\; \hat{\tau}_{\text{err}}\Big),
\label{eq:verification_tuple}
\end{equation}
where $\hat{y}_{i,k} \in \mathcal{Y}$ (comprising \textsc{Supported}, \textsc{Partially Supported}, \textsc{Contradicted}, and \textsc{Neutral}) is the point prediction, $\mathcal{C}_{\alpha} \subseteq \mathcal{Y}$ is a conformal prediction set satisfying $\mathbb{P}(y^*_{i,k} \in \mathcal{C}_{\alpha}) \geq 1-\alpha$, $\Delta_{\mathcal{E}} \in [-1, 1]$ quantifies the epistemic modality shift, and $\hat{\tau}_{\text{err}}$ is the diagnosed misinterpretation category.

\subsection{Research Questions and Measurable Objectives}
To address the identified gaps, this study investigates four formal research questions:
\begin{itemize}
\item \textbf{RQ1 (Discourse-Scoped Bipartite Claim Attribution):} How can multi-sentence discourse contexts and compound citing sentences containing multiple inline references be mathematically decomposed into decontextualized, citation-specific atomic propositions without cross-reference interference?
\item \textbf{RQ2 (Structure-Aware Multimodal Evidence Retrieval):} How can lexical, dense citation-informed, late-interaction, and tabular/visual representations be fused with rhetorical section priors to maximize evidence recall across full-text manuscripts?
\item \textbf{RQ3 (Epistemic Modality Delta and Scope Alignment):} How can continuous shifts in epistemic commitment (hedging erasure, causal inflation), experimental scope, and numerical values be formally quantified beyond discrete tri-state NLI?
\item \textbf{RQ4 (Risk-Controlled Conformal Auditing):} How can split conformal prediction calibrate verification uncertainty so that automated flags maintain a guaranteed upper bound $\alpha$ on false coverage errors while routing ambiguous claims to human reviewers?
\end{itemize}

Aligned with RQ1--RQ4, our measurable research objectives are: (O1) to design a cohesion-penalized window selector and bipartite clause-to-citation mapping algorithm; (O2) to formulate a section-weighted multimodal Reciprocal Rank Fusion (RRF) and cross-encoder reranking pipeline; (O3) to derive a composite verification score integrating NLI entailment, Epistemic Modality Delta $\Delta_{\mathcal{E}}$, scope divergence $D_{\text{scope}}$, and numerical tolerance $\delta_{\text{num}}$; and (O4) to construct a split conformal prediction wrapper and reproducible evaluation protocol.

\begin{table*}[t]
\centering
\footnotesize
\caption{Analytical Taxonomy of Misinterpreted but Genuine Citations Grounded in Empirical Literature}
\label{tab:misinterpretation_taxonomy}
\renewcommand{\arraystretch}{1.14}
\begin{tabularx}{\textwidth}{L{0.5cm} L{3.0cm} L{4.2cm} L{4.2cm} Y}
\toprule
\textbf{ID} & \textbf{Misinterpretation Class} & \textbf{Formal Epistemic Definition} & \textbf{Linguistic / Structural Signature} & \textbf{Primary Detection Signal in CEDAR} \\
\midrule
T1 & Speculative Elevation (Hedging Erasure) & Source presents an unverified hypothesis or future direction; citing text asserts it as proven empirical fact \cite{greenberg2009citation,hyland1998hedging}. & Deletion of modal auxiliaries (\emph{may, might, could}) and epistemic verbs (\emph{suggests, hypothesizes}); source located in Intro/Discussion. & High positive Epistemic Delta ($\Delta_{\mathcal{E}} > \tau_{\text{mod}}$) and low section authority weight $w_{z(e)}$. \\
T2 & Causal Inflation & Source reports observational correlation or statistical association; citing text asserts a direct causal mechanism \cite{jergas2015quotation,mogull2017accuracy}. & Substitution of correlational predicates (\emph{associated with, linked to}) with causal verbs (\emph{causes, induces, drives}) despite observational study design. & Predicate-role mismatch between study design metadata and causal verb class in $c_k$. \\
T3 & Scope \& Population Extrapolation & Source evaluates a restricted cohort (e.g., murine \emph{in vitro} cells or synthetic data); citing text generalizes to universal/clinical populations \cite{begley2012raise,pavlovic2021accurate}. & Removal of domain restrictors or insertion of universal quantifiers (\emph{all, universally, clinically proven in patients}). & Entity/population divergence $D_{\text{scope}}(c_k, E^*) > \tau_{\text{scope}}$ across PICO/benchmark slots. \\
T4 & Directional Inversion (Finding Reversal) & Source reports a null effect, statistically non-significant trend, or opposite outcome; citing text claims positive support \cite{delacey1985accurate,todd2010marine}. & Polarity inversion or neglect of negation markers (\emph{no significant difference, failed to improve, $p > 0.05$}) in source Results. & High cross-encoder \textsc{Contradicted} posterior $P_{\Theta}(\textsc{Contradicted} \mid c_k, E^*)$. \\
T5 & Quantitative \& Statistical Distortion & Direct misreporting of effect sizes, percentages, sample sizes $n$, $p$-values, or confidence intervals \cite{harzing2002referencing,rekdal2014urban}. & Order-of-magnitude decimal shifts, metric swapping, or reporting an upper confidence bound as a mean point estimate. & Neuro-symbolic numerical inequality/tolerance check $\delta_{\text{num}}(c_k, E^*) > \epsilon_{\text{tol}}$ over prose and tables \cite{lu2023scitab}. \\
T6 & Partial \& Compound Conflation & Citing sentence combines multiple sub-claims where the source supports one clause but is silent on or refutes the others \cite{kamoi2023wice,min2023factscore}. & Multi-clause coordination or subordination anchored to a single citation or mismatched multi-citation brackets. & Sub-claim entailment variance across decomposed atomic units $\{c_{i,1}, \dots, c_{i,m_i}\}$. \\
T7 & Caveat \& Boundary Erasure & Source explicitly conditions a finding on strict assumptions or limitations; citing text states the result unconditionally \cite{hyland1998hedging,smith2020quotation}. & Omission of conditional clauses (\emph{under alkaline conditions, assuming IID noise, when $d < 10$}). & Discourse condition-slot drop detected between source rationale $E^*$ and atomic claim $c_k$. \\
T8 & Secondary Citation Cascade (Straw-Man / Herd) & Citing author copies a distorted claim from an intermediate secondary paper without consulting the primary work \cite{simkin2003read,rekdal2014sinking}. & Identical phrasing or shared numerical/page misprints matching an intermediate citing paper rather than $D_{\text{src}}^{(j)}$. & Zero lexical/semantic support in $D_{\text{src}}^{(j)}$ paired with high similarity to intermediate citing nodes. \\
\bottomrule
\end{tabularx}
\end{table*}

\section{Taxonomy and Case Analysis of Citation Misinterpretation}
\label{sec:taxonomy}

\subsection{Analytical Taxonomy of Citation Misinterpretation}
Synthesizing empirical findings from medical, biological, and social science quotation audits \cite{delacey1985accurate,greenberg2009citation,harzing2002referencing,jergas2015quotation,mogull2017accuracy,pavlovic2021accurate,smith2020quotation}, we organize misinterpreted but genuine citations into eight operational classes (T1--T8), defined in Table~\ref{tab:misinterpretation_taxonomy}. Unlike coarse binary error definitions, this taxonomy separates rhetorical modality distortions (T1, T2, T7) from domain-transfer errors (T3), polarity reversals (T4), quantitative discrepancies (T5), syntactic decomposition failures (T6), and network propagation effects (T8).

\subsection{Authentic Historical Cases Versus Illustrative Diagnostic Scenarios}
A critical methodological requirement in citation research is distinguishing between \emph{authentic, literature-documented historical cases} of citation distortion and \emph{illustrative diagnostic scenarios} constructed to test or explain algorithmic components \cite{rekdal2014sinking,rekdal2014urban}. Below, we first examine three verified historical cases from the peer-reviewed scientometric literature, followed by three synthetic diagnostic scenarios that illustrate how CEDAR's mathematical modules process domain-specific failure modes.

\subsubsection{Verified Historical Cases from Published Literature:}
\begin{itemize}
\item \textbf{Historical Case 1 ($\beta$-Amyloid in Inclusion Body Myositis --- Greenberg \cite{greenberg2009citation}):} Greenberg analyzed 242 published articles and 675 citations examining whether $\beta$-amyloid precursor protein is pathologically expressed in skeletal muscle fibers in inclusion body myositis (IBM). An early paper offered a tentative hypothesis in its discussion section, which subsequent papers cited as an established experimental finding (Class T1: Speculative Elevation). Meanwhile, primary experimental studies that failed to find $\beta$-amyloid mRNA differences in IBM muscle were either uncited or cited in ways that obscured their negative results (Class T4).
\item \textbf{Historical Case 2 (Expatriate Failure Rates --- Harzing \cite{harzing2002referencing}):} Harzing traced widespread claims in international human resource management literature asserting that 16\% to 40\% of American expatriate managers fail and return prematurely from foreign assignments. Auditing the citation trail back to the primary sources revealed that the original empirical study did not measure actual expatriate failure rates at all, but rather surveyed personnel managers' rough perceptions, which subsequent authors repeatedly copied and elevated into hard statistical facts (Classes T1, T5, and T8).
\item \textbf{Historical Case 3 (Decimal Shift and Secondary Copying --- Rekdal \cite{rekdal2014urban}):} Rekdal documented how quantitative errors and unverified secondary citations propagate across decades of scientific literature as ``academic urban legends,'' demonstrating how a tenfold decimal misplacement or a speculative remark in a secondary source becomes canonized when authors cite primary sources indirectly without consulting the original text (Classes T5 and T8).
\end{itemize}

\subsubsection{Illustrative Diagnostic Scenarios for Pipeline Walkthrough:}
To illustrate how an automated verification pipeline must process fine-grained textual, tabular, and numerical divergences across disciplines, we formulate three synthetic diagnostic scenarios (explicitly designated as illustrative test cases rather than published empirical audits):
\begin{itemize}
\item \textbf{Diagnostic Scenario A (Biomedical Model-to-Clinic Extrapolation --- Class T3):} Suppose a citing manuscript asserts: \emph{``Inhibiting pathway P38 via compound X clinically reverses cognitive decline in late-stage human patients.''} If the cited target paper $D_{\text{src}}$ actually reports an \emph{in vitro} neuroblastoma assay and a 14-day transgenic mouse trial ($n=8$) while explicitly cautioning in its Discussion that human clinical efficacy remains untested, a surface-level bi-encoder retrieves high topical similarity, whereas CEDAR flags a severe population/scope divergence $D_{\text{scope}}$ and caveat erasure (Class T3/T7).
\item \textbf{Diagnostic Scenario B (Machine Learning Benchmark Overgeneralization --- Class T3/T6):} Suppose a citing sentence states: \emph{``Sparse Mixture-of-Experts (MoE) transformers consistently outperform dense models across all downstream NLP reasoning and translation benchmarks.''} If the cited paper's prose reports gains on pre-training perplexity and high-resource translation, but its ablation table (e.g., Table~4) shows that dense models outperform the MoE variant on multi-step mathematical reasoning and low-resource translation, verifying this citation requires joint reasoning over narrative text and structured tables \cite{herzig2020tapas,lu2023scitab} to refute the universal quantifier \emph{``across all benchmarks.''}
\item \textbf{Diagnostic Scenario C (Quantitative Decimal Shift in Materials Science --- Class T5):} Suppose a citing sentence claims: \emph{``Nitrogen-doped graphene achieves a 45\% enhancement in peak current density.''} If the cited source's results section and figure caption state a \emph{4.5\% increase} under alkaline conditions, lexical and dense embedding models assign near-identical representations to ``45\%'' and ``4.5\%,'' whereas CEDAR's neuro-symbolic numerical verification module detects a tenfold relative discrepancy ($\delta_{\text{num}} = |45 - 4.5|/4.5 = 9.0 \gg \epsilon_{\text{tol}}$) and stripped boundary conditions (Classes T5 and T7).
\end{itemize}

\begin{figure*}[t]
\centering
\includegraphics[width=0.96\textwidth]{images/fig1_system_architecture.pdf}
\caption{End-to-end architecture of the proposed 12-stage CEDAR framework for detecting misinterpreted but genuine citations. Tier~1 ingests compiled PDFs, optimizes discourse context windows, and executes bipartite atomic claim-to-anchor decomposition. Tier~2 verifies bibliographic authenticity and performs section-weighted multimodal retrieval across prose, tables, and figure captions. Tier~3 computes NLI posteriors, Epistemic Modality Delta ($\Delta_{\mathcal{E}}$), numerical consistency, and split conformal prediction sets ($\mathcal{C}_{\alpha}$) to generate explainable human-in-the-loop Verification Cards.}
\label{fig:architecture}
\end{figure*}

\section{Proposed CEDAR Verification Architecture}
\label{sec:framework}

To overcome the limitations of isolated sentence matching and abstract-only verification, CEDAR structures automated citation auditing into a 12-stage modular pipeline organized across three functional tiers, depicted in Fig.~\ref{fig:architecture} and summarized systematically in Table~\ref{tab:pipeline_stages}.

\subsection{Tier 1: Document Ingestion, Discourse Scoping, and Bipartite Decomposition}
\subsubsection{Stage 1 (Document Layout Parsing and Anchor Isolation):}
\emph{Input:} Raw compiled PDF of the citing manuscript $D_{\text{cite}}$. \emph{Processing:} The ingestion module applies layout-aware sequence labeling via GROBID \cite{lopez2009grobid} to convert multi-column PDF streams into structured Text Encoding Initiative (TEI) XML, isolating sentence boundaries $\mathcal{S} = (s_1, \dots, s_M)$, inline citation anchors $a_i \in \mathcal{A}$ (handling both numbered brackets $[j]$ and author--year strings), and structured bibliography records $\mathcal{B}$. \emph{Output:} Indexed sentence sequence $\mathcal{S}$ and citation linkage map $\mu(a_i) = b_j$. \emph{Justification \& Limitations:} Accurate XML tree extraction prevents footnote, header, and table-caption contamination, though heavily corrupted legacy PDF scans can still introduce OCR tokenization noise.

\subsubsection{Stage 2 (Cohesion-Weighted Discourse Window Extraction):}
\emph{Input:} Focal citing sentence $s_{t(i)}$ and surrounding paragraph sentences $(s_{t-W}, \dots, s_{t+W})$. \emph{Processing:} Rather than using a fixed single sentence ($w=0$), which drops premises introduced in prior sentences \cite{jurgens2018measuring}, CEDAR evaluates discourse cohesion ties (anaphoric pronouns such as \emph{this method, their study}, lexical entity overlap, and rhetorical connectives such as \emph{consequently, in contrast}) to select an asymmetric window $C(a_i) = (s_{t-w^-}, \dots, s_t, \dots, s_{t+w^+})$. \emph{Output:} Discourse context window $C(a_i)$. \emph{Justification \& Limitations:} Captures multi-sentence arguments concluding with a citation; fails only when an author places a citation at the end of a paragraph to cover an entire multi-page section without cohesive markers.

\begin{figure}[t]
\centering
\includegraphics[width=\columnwidth]{images/fig2_claim_decomposition.pdf}
\caption{Discourse-scoped atomic claim decomposition and bipartite citation-anchor attribution (Stages 2--4). Resolving pronominal references and splitting subordinate contrast clauses prevents cross-reference contamination when verifying multi-citation sentences.}
\label{fig:decomposition}
\end{figure}

\subsubsection{Stage 3 (Coreference Resolution and Atomic Claim Decomposition):}
\emph{Input:} Discourse context $C(a_i)$. \emph{Processing:} Following atomic decomposition principles established in WiCE \cite{kamoi2023wice} and FActScore \cite{min2023factscore}, an instruction-tuned scientific language model constrained by dependency parse trees resolves pronominal anaphora to explicit canonical entities and splits compound/subordinate sentences into minimal, self-contained propositions $\{c_1, \dots, c_m\}$. \emph{Output:} Set of decontextualized atomic claims. \emph{Justification \& Limitations:} Eliminates false negatives caused by unresolved pronouns (e.g., ``it'') during passage retrieval; over-decomposition of tightly coupled conditional clauses must be prevented by preserving ``if/when'' condition modifiers inside each atomic proposition.

\subsubsection{Stage 4 (Bipartite Claim-to-Citation Attribution):}
\emph{Input:} Atomic propositions $\{c_1, \dots, c_m\}$, dependency tree of $s_{t(i)}$, and inline anchors $\{a_1, \dots, a_r\}$ in $C(a_i)$. \emph{Processing:} As illustrated in Fig.~\ref{fig:decomposition}, when a sentence contains multiple citations (e.g., contrasting Model~A $[12]$ with Model~B $[18]$), a syntactic scope resolver computes dependency-path distance and clause-boundary dominance to construct a bipartite graph matching each atomic claim $c_k$ exclusively to its governing citation anchor(s) $a_i$. \emph{Output:} Attributed claim--source tuples $(c_{i,k}, b_j)$. \emph{Justification \& Limitations:} Prevents verifying a clause meant for Reference~$[18]$ against Reference~$[12]$; remains ambiguous when authors place a bundled bracket $[12\text{--}16]$ at the end of a synthesis sentence without per-clause attribution.

\subsection{Tier 2: Bibliographic Resolution and Multimodal Hierarchical Retrieval}
\subsubsection{Stage 5 (Bibliographic Authenticity and Horizon Verification):}
\emph{Input:} Bibliography entry $b_j = \mu(a_i)$. \emph{Processing:} Queries Crossref, OpenAlex, and PubMed APIs to verify that $b_j$ resolves to an authentic publication $D_{\text{src}}^{(j)}$ and checks title/author/year consistency. If full text is accessible via OpenAlex, PubMed Central, or S2ORC \cite{lo2020s2orc}, the epistemic horizon flag is set to $h_{\text{doc}} = 1$; if commercial paywalls \cite{lariviere2015oligopoly} restrict access to the abstract only, $h_{\text{doc}} = 0$ is logged. \emph{Output:} Verified source document $D_{\text{src}}^{(j)}$ and horizon indicator $h_{\text{doc}} \in \{0, 1\}$. \emph{Justification \& Limitations:} Separates phantom/hallucinated references from genuine publications upfront and prevents false \textsc{Contradicted} verdicts when full text is behind a paywall.

\subsubsection{Stage 6 (Multimodal Source Document Indexing):}
\emph{Input:} Full-text PDF or XML of $D_{\text{src}}^{(j)}$. \emph{Processing:} Segments narrative text into sliding 150-word passages labeled by rhetorical zone $z(e_r)$ \cite{teufel2006automatic}, while extracting structured tables into row-column linearized representations compatible with TaPas \cite{herzig2020tapas} and SCITAB \cite{lu2023scitab} and extracting figure captions alongside chart-derived data series \cite{akhtar2023reading}. \emph{Output:} Unified multimodal evidence index $D_{\text{src}}^{(j)} = \{e_1, \dots, e_{N_j}\}$. \emph{Justification \& Limitations:} Ensures that ablation metrics in tables and trends in figure captions are retrievable as first-class evidentiary candidates.

\subsubsection{Stage 7 (Section-Weighted Hybrid Reciprocal Rank Fusion):}
\emph{Input:} Atomic claim $c_{i,k}$ and source evidence pool $\{e_r\}_{r=1}^{N_j}$. \emph{Processing:} Executes three parallel retrievers: (i) sparse BM25 lexical scoring \cite{robertson2009probabilistic} for exact entity/gene/metric matching; (ii) dense bi-encoder cosine similarity using domain-adapted scientific encoders (SciBERT \cite{beltagy2019scibert}, SBERT \cite{reimers2019sentence}, DPR \cite{karpukhin2020dense}, SPECTER \cite{cohan2020specter}); and (iii) token-level late interaction via ColBERT \cite{khattab2020colbert}. Rankings are combined via Reciprocal Rank Fusion (RRF) \cite{cormack2009reciprocal} weighted by rhetorical section priors $w_{z(e_r)}$ that prioritize \emph{Results}, \emph{Tables}, and \emph{Methods} over \emph{Introduction} background remarks. \emph{Output:} Top-$K_1$ candidate evidentiary units ($K_1 = 25$). \emph{Justification \& Limitations:} Overcomes the ``topical overlap trap'' of pure bi-encoders while retaining sub-second retrieval latency.

\subsubsection{Stage 8 (Cross-Encoder Fine-Grained Reranking):}
\emph{Input:} Claim $c_{i,k}$ and top-$K_1$ candidate units. \emph{Processing:} Concatenates $\texttt{[CLS]} \oplus c_{i,k} \oplus \texttt{[SEP]} \oplus e_r$ through a deep cross-attention transformer backbone to perform all-to-all token interaction and select the minimal rationale set $E^*(c_{i,k})$ of size $K_2 \leq 5$. \emph{Output:} Ranked evidentiary rationale set $E^*$. \emph{Justification \& Limitations:} Captures fine-grained negation, comparatives, and modifiers before invoking expensive reasoning modules.

\begin{table*}[t]
\centering
\footnotesize
\caption{Stage-by-Stage Specification of the Proposed 12-Stage CEDAR Verification Pipeline}
\label{tab:pipeline_stages}
\renewcommand{\arraystretch}{1.13}
\begin{tabularx}{\textwidth}{C{0.5cm} L{2.5cm} L{2.3cm} L{2.4cm} L{3.2cm} Y}
\toprule
\textbf{Stg.} & \textbf{Module Name} & \textbf{Primary Input} & \textbf{Primary Output} & \textbf{Core Algorithm / Model} & \textbf{Technical Justification \& Known Limitation} \\
\midrule
1 & Layout \& XML Ingestion & Raw citing PDF $D_{\text{cite}}$ & TEI-XML sentences $\mathcal{S}$, anchors $\mathcal{A}$, bib $\mathcal{B}$ & GROBID \cite{lopez2009grobid} CRF/BiLSTM sequence parsers & Isolates prose from headers/footnotes; degraded by low-DPI scanned legacy PDFs. \\
2 & Discourse Window Selection & Focal sentence $s_{t(i)}$ & Context window $C(a_i)$ (Eq.~\ref{eq:window_opt}) & Cohesion-penalized boundary optimization \cite{jurgens2018measuring} & Recovers premises across sentence boundaries; cannot resolve unanchored whole-section citations. \\
3 & Atomic Claim Decomposition & Context window $C(a_i)$ & Decontextualized claims $\{c_1, \dots, c_m\}$ & Dependency-guided LLM decomposition \cite{kamoi2023wice,min2023factscore} & Prevents pronoun ambiguity and multi-clause entanglement; requires condition preservation. \\
4 & Bipartite Anchor Attribution & Claims $\{c_k\}$ \& anchors $\{a_i\}$ & Attributed pairs $(c_{i,k}, b_j)$ (Eq.~\ref{eq:bipartite}) & Syntactic dependency-tree dominance matching & Eliminates multi-reference cross-talk; bundled brackets $[1\text{--}5]$ require exhaustive check. \\
5 & Bibliographic \& Horizon Check & Bib entry $b_j = \mu(a_i)$ & Resolved $D_{\text{src}}^{(j)}$ \& flag $h_{\text{doc}} \in \{0,1\}$ & Crossref/OpenAlex/S2ORC API linkage \cite{lo2020s2orc} & Filters hallucinated references and logs paywall truncation ($h_{\text{doc}}=0$) \cite{lariviere2015oligopoly}. \\
6 & Multimodal Indexing & Full text of $D_{\text{src}}^{(j)}$ & Units $\mathcal{U}_{\text{prose}} \cup \mathcal{U}_{\text{table}} \cup \mathcal{U}_{\text{fig}}$ & Rhetorical zoning \cite{teufel2006automatic}, TaPas \cite{herzig2020tapas}, SCITAB \cite{lu2023scitab} & Exposes numerical tables and figure captions to retrieval; complex multi-panel blots remain hard. \\
7 & Section-Weighted Hybrid RRF & Claim $c_{i,k}$ \& pool $\{e_r\}$ & Top-$K_1$ candidate passages ($K_1=25$) & BM25 \cite{robertson2009probabilistic} + SPECTER \cite{cohan2020specter} + ColBERT \cite{khattab2020colbert} RRF \cite{cormack2009reciprocal} & Combines exact entity matching with semantic recall and section authority priors $w_{z(e)}$. \\
8 & Cross-Encoder Reranking & Claim $c_{i,k}$ \& top-$K_1$ units & Minimal rationale set $E^*$ ($K_2 \leq 5$) & Cross-attention transformer scorer \cite{wadden2022multivers} & Resolves bi-encoder topical overlap trap via token-by-token cross-attention. \\
9 & Scientific NLI Adjudication & Pair $(c_{i,k}, E^*)$ & Probability vector $\hat{\mathbf{p}}(y \mid c_{i,k}, E^*)$ & MultiVerS / SciFact-tuned NLI classifier \cite{wadden2020fact,wadden2022multivers} & Evaluates multi-paragraph propositional entailment; restricted by context window length. \\
10 & Epistemic Delta \& Num.\ Check & Pair $(c_{i,k}, E^*)$ & $\Delta_{\mathcal{E}}, D_{\text{scope}}, \delta_{\text{num}}$ (Eqs.~\ref{eq:epistemic_delta}--\ref{eq:num_check}) & Modality scorer \cite{hyland1998hedging} + symbolic quantity parser \cite{lu2023scitab} & Detects hedging erasure, causal inflation, and order-of-magnitude decimal shifts. \\
11 & Split Conformal Calibration & Posteriors $\hat{\mathbf{p}}$ \& calib.\ set $\mathcal{D}_{\text{cal}}$ & Prediction set $\mathcal{C}_{\alpha}(c_{i,k}, E^*)$ (Eq.~\ref{eq:conformal_set}) & Split conformal set-valued classification \cite{sadinle2019least,angelopoulos2023conformal} & Guarantees $1-\alpha$ coverage and abstains under epistemic ambiguity ($|\mathcal{C}_{\alpha}| \geq 2$). \\
12 & Explainable Verification Card & All stage outputs & Structured JSON/UI Audit Card & Token-level rationales \cite{rashkin2023measuring,gao2023enabling} \& severity $\mathcal{S}_{\text{sev}}$ & Provides transparent audit trail for editors and authors; keeps final judgment human-in-the-loop. \\
\bottomrule
\end{tabularx}
\end{table*}

\subsection{Tier 3: Epistemic Inference, Calibration, and Audit Generation}
\subsubsection{Stage 9 (Deep Scientific NLI Adjudication):}
\emph{Input:} Atomic claim $c_{i,k}$ and retrieved rationale set $E^*(c_{i,k})$. \emph{Processing:} Computes the posterior probability distribution $\hat{\mathbf{p}}(y \mid c_{i,k}, E^*)$ over $y \in \mathcal{Y}$ using a scientific cross-encoder trained with multi-task rationale supervision \cite{wadden2020fact,wadden2022multivers}.

\subsubsection{Stage 10 (Epistemic Modality Delta, Scope, and Numerical Verification):}
\emph{Input:} Claim $c_{i,k}$ and rationale $E^*$. \emph{Processing:} Evaluates three specialized consistency checks: (i) the Epistemic Modality Delta $\Delta_{\mathcal{E}}$ comparing linguistic commitment markers \cite{hyland1998hedging} in $c_{i,k}$ against $E^*$; (ii) structured PICO/benchmark scope alignment $D_{\text{scope}}$; and (iii) neuro-symbolic numerical and statistical inequality checking $\delta_{\text{num}}$ over extracted quantities, units, and $p$-values \cite{lu2023scitab,akhtar2023reading}.

\subsubsection{Stage 11 (Split Conformal Uncertainty Calibration):}
\emph{Input:} Calibrated posterior probabilities and non-conformity scores over a held-out calibration corpus $\mathcal{D}_{\text{cal}}$. \emph{Processing:} Constructs a distribution-free prediction set $\mathcal{C}_{\alpha}(c_{i,k}, E^*) \subseteq \mathcal{Y}$ guaranteed to contain the true epistemic relationship with probability at least $1-\alpha$ \cite{sadinle2019least,angelopoulos2023conformal}, as shown in Fig.~\ref{fig:calibration}. Singleton sets ($|\mathcal{C}_{\alpha}| = 1$) yield high-confidence automated classifications, whereas multi-label sets ($|\mathcal{C}_{\alpha}| \geq 2$) trigger automatic abstention and human review.

\subsubsection{Stage 12 (Explainable Verification Card and Human-in-the-Loop Audit):}
\emph{Input:} Complete verification tuple $\mathcal{V}(c_{i,k}, D_{\text{src}}^{(j)})$. \emph{Processing:} Generates an interactive Verification Card highlighting the exact spans in $c_{i,k}$ responsible for positive $\Delta_{\mathcal{E}}$ or numerical divergence alongside verbatim quotes and table coordinates from $D_{\text{src}}^{(j)}$ \cite{rashkin2023measuring,gao2023enabling}.

\begin{figure*}[t]
\centering
\includegraphics[width=0.95\textwidth]{images/fig3_epistemic_calibration.pdf}
\caption{Epistemic Modality Delta ($\Delta_{\mathcal{E}}$) and Split Conformal Risk Control in CEDAR (Stages 10--11). (a) Mapping citing claim strength $S_{\text{claim}}(c_k)$ against source evidence strength $S_{\text{evid}}(E^*)$ separates acceptable hedging ($|\Delta_{\mathcal{E}}| \leq \tau_{\text{mod}}$) from speculative elevation and causal inflation ($\Delta_{\mathcal{E}} > \tau_{\text{mod}}$). (b) Split conformal calibration thresholds the non-conformity score $s(c_k, E^*, y)$ at quantile $\hat{q}_{1-\alpha}$, issuing autonomous flags only for unambiguous singleton sets ($|\mathcal{C}_{\alpha}|=1$) and routing ambiguous sets ($|\mathcal{C}_{\alpha}| \geq 2$) to human auditors.}
\label{fig:calibration}
\end{figure*}

\section{Mathematical Formulation and Algorithmic Design}
\label{sec:math}

We now formalize the mathematical equations governing each core computation in the CEDAR pipeline, defining every variable explicitly.

\subsection{Cohesion-Penalized Discourse Window Selection}
Let $s_t$ be the focal sentence containing citation anchor $a_i$, and let $W_{\max}$ be the maximum permissible window radius within the current paragraph (typically $W_{\max}=3$). For candidate backward and forward offsets $(w^-, w^+) \in \{0, \dots, W_{\max}\}^2$, let $C_{w^-, w^+}(a_i) = (s_{t-w^-}, \dots, s_t, \dots, s_{t+w^+})$. We define the optimal discourse window $(w^{-*}, w^{+*})$ by maximizing a cohesion-minus-drift objective:
\begin{equation}
\begin{aligned}
(w^{-*}, w^{+*}) = \arg\max_{0 \leq w^-, w^+ \leq W_{\max}} \Bigg[ & \sum_{u=t-w^-}^{t+w^+} \gamma^{|u-t|} \, \Phi_{\text{coh}}(s_u, s_t) \\
& - \lambda_{\text{len}} (w^- + w^+) \\
& - \lambda_{\text{cite}} \sum_{u \neq t} \mathbb{I}(\mathcal{A}_u \setminus \mathcal{A}_t \neq \emptyset) \Bigg],
\end{aligned}
\label{eq:window_opt}
\end{equation}
where $\gamma \in (0, 1)$ is a discourse distance decay factor, $\lambda_{\text{len}} > 0$ penalizes unnecessary window expansion, $\mathcal{A}_u$ denotes the set of citation anchors appearing in sentence $s_u$, $\mathbb{I}(\cdot)$ is the indicator function, and $\lambda_{\text{cite}} > 0$ penalizes absorbing adjacent sentences anchored to unrelated citations. The pairwise discourse cohesion function $\Phi_{\text{coh}}(s_u, s_t) \in [0, 1]$ is defined as:
\begin{equation}
\begin{aligned}
\Phi_{\text{coh}}(s_u, s_t) = \;& \beta_1 \psi_{\text{coref}}(s_u, s_t) + \beta_2 \psi_{\text{conn}}(s_u, s_t) \\
& + \beta_3 \frac{\langle \mathbf{h}(s_u), \mathbf{h}(s_t) \rangle}{\|\mathbf{h}(s_u)\|_2 \|\mathbf{h}(s_t)\|_2},
\end{aligned}
\label{eq:cohesion_score}
\end{equation}
where $\psi_{\text{coref}}(s_u, s_t) \in \{0, 1\}$ indicates whether $s_t$ contains an unresolved anaphoric expression whose antecedent lies in $s_u$, $\psi_{\text{conn}}(s_u, s_t) \in [0, 1]$ measures discourse connective continuity, $\mathbf{h}(\cdot) \in \mathbb{R}^d$ is a sentence embedding \cite{beltagy2019scibert,reimers2019sentence}, and $\beta_1 + \beta_2 + \beta_3 = 1$.

\subsection{Bipartite Claim Decomposition and Citation-Anchor Mapping}
Let $\mathcal{C}(C^*(a_i)) = \{c_1, c_2, \dots, c_m\}$ denote the set of decontextualized atomic propositions extracted from the optimal window $C^*(a_i)$, and let $\mathcal{A}^* = \{a_1, \dots, a_r\}$ denote all inline citation anchors present in $C^*(a_i)$. Let $\ell(c_k, a_i)$ denote the syntactic affinity logit between atomic claim $c_k$ and anchor $a_i$:
\begin{equation}
\begin{aligned}
\ell(c_k, a_i) = \;& -\eta_1 d_{\text{dep}}(\pi(c_k), a_i) - \eta_2 d_{\text{tok}}(\pi(c_k), a_i) \\
& + \eta_3 \mathbb{I}_{\text{clause}}(\pi(c_k), a_i),
\end{aligned}
\label{eq:bipartite_logit}
\end{equation}
where $\pi(c_k)$ is the syntactic head predicate of the source clause from which $c_k$ was extracted, $d_{\text{dep}}(\pi(c_k), a_i)$ is the shortest path length between $\pi(c_k)$ and anchor $a_i$ in the sentence dependency graph, $d_{\text{tok}}(\pi(c_k), a_i)$ is linear token distance, $\mathbb{I}_{\text{clause}}(\pi(c_k), a_i) \in \{0, 1\}$ equals $1$ if $a_i$ is governed by the minimal finite clause dominating $\pi(c_k)$, and $\eta_1, \eta_2, \eta_3 > 0$ are positive weights. The normalized attribution probability $A(c_k, a_i)$ is then:
\begin{equation}
A(c_k, a_i) = \frac{\exp(\ell(c_k, a_i))}{\sum_{a_{i'} \in \mathcal{A}^*} \exp(\ell(c_k, a_{i'}))}.
\label{eq:bipartite}
\end{equation}
Atomic claim $c_k$ is attributed to anchor $a_i$ if and only if $A(c_k, a_i) \geq \tau_{\text{attr}}$ or $a_i = \arg\max_{a_{i'} \in \mathcal{A}^*} A(c_k, a_{i'})$, where $\tau_{\text{attr}} \in (0, 1)$ is the multi-anchor attribution threshold.

\subsection{Section-Weighted Multimodal Reciprocal Rank Fusion}
Given an attributed atomic claim $c_k$ and resolved source document $D_{\text{src}}^{(j)} = \{e_1, \dots, e_{N_j}\}$, each evidentiary unit $e_r$ (whether a narrative paragraph, a linearized table row/column from TaPas/SCITAB \cite{herzig2020tapas,lu2023scitab}, or a figure caption \cite{akhtar2023reading}) is scored by three complementary retrieval functions:
\begin{enumerate}
\item \emph{Sparse Lexical Matching (BM25 \cite{robertson2009probabilistic}):}
\begin{equation}
S_{\text{lex}}(c_k, e_r) = \sum_{q \in c_k} \text{IDF}(q) \cdot \frac{f(q, e_r)\,(k_1 + 1)}{f(q, e_r) + k_1 \left(1 - b + b \frac{|e_r|}{L_{\text{avg}}}\right)},
\label{eq:bm25}
\end{equation}
where $f(q, e_r)$ is the term frequency of token $q$ in unit $e_r$, $|e_r|$ is the token length of $e_r$, $L_{\text{avg}}$ is the mean unit length across $D_{\text{src}}^{(j)}$, and $k_1, b$ are standard BM25 saturation and length-normalization parameters.
\item \emph{Dense Bi-Encoder Similarity \cite{beltagy2019scibert,cohan2020specter,karpukhin2020dense}:}
\begin{equation}
S_{\text{dense}}(c_k, e_r) = \frac{\langle E_{\text{claim}}(c_k), E_{\text{evid}}(e_r) \rangle}{\|E_{\text{claim}}(c_k)\|_2 \, \|E_{\text{evid}}(e_r)\|_2},
\label{eq:dense_sim}
\end{equation}
where $E_{\text{claim}}$ and $E_{\text{evid}}$ are domain-adapted transformer encoders.
\item \emph{Contextualized Late Interaction (ColBERT MaxSim \cite{khattab2020colbert}):}
\begin{equation}
S_{\text{late}}(c_k, e_r) = \sum_{u=1}^{|c_k|} \max_{v=1, \dots, |e_r|} \langle \mathbf{E}_c^{(u)}, \mathbf{E}_e^{(v)} \rangle,
\label{eq:colbert}
\end{equation}
where $\mathbf{E}_c^{(u)}, \mathbf{E}_e^{(v)} \in \mathbb{S}^{d-1}$ denote $\ell_2$-normalized contextualized token embeddings for the $u$-th token of $c_k$ and $v$-th token of $e_r$.
\end{enumerate}

Let $\text{rank}_m(c_k, e_r) \in \{1, \dots, N_j\}$ denote the rank assigned to unit $e_r$ by retrieval modality $m \in \mathcal{M} = \{\text{lex}, \text{dense}, \text{late}\}$. To incorporate rhetorical section authority (preventing speculative remarks in an \emph{Introduction} from outranking empirical data in \emph{Results} or \emph{Tables}), we define the Section-Weighted Reciprocal Rank Fusion (SW-RRF) score by extending Cormack et al.\ \cite{cormack2009reciprocal}:
\begin{equation}
S_{\text{SW-RRF}}(c_k, e_r) = w_{z(e_r)} \sum_{m \in \{\text{lex}, \text{dense}, \text{late}\}} \frac{\omega_m}{\kappa_{\text{rrf}} + \text{rank}_m(c_k, e_r)},
\label{eq:sw_rrf}
\end{equation}
where $\kappa_{\text{rrf}} = 60$ is the RRF smoothing constant, $\omega_m > 0$ ($\sum_m \omega_m = 1$) weights each retrieval view, and $w_{z(e_r)} \in (0, 1]$ is the rhetorical section prior (with $w_{\text{Results}} = w_{\text{Table}} = 1.0$, $w_{\text{Methods}} = w_{\text{Caption}} = 0.9$, $w_{\text{Abstract}} = 0.85$, $w_{\text{Discussion}} = 0.75$, and $w_{\text{Intro}} = 0.50$). The top-$K_1$ candidates under $S_{\text{SW-RRF}}$ are then reranked by a cross-attention encoder $g_{\text{cross}}(c_k, e_r)$ to select the top-$K_2$ evidentiary rationale set $E^*(c_k) = \{e_{r_1}, \dots, e_{r_{K_2}}\}$.

\subsection{Epistemic Modality Delta, Scope Divergence, and Numerical Verification}
Standard NLI models output a categorical probability vector $\hat{\mathbf{p}}_{\Theta}(y \mid c_k, E^*)$ over $\mathcal{Y}$, yet fail when $c_k$ and $E^*$ share topical entities while diverging in epistemic certainty, experimental scope, or numerical magnitude. We formalize three complementary verification metrics:

\subsubsection{Epistemic Modality Delta (\texorpdfstring{$\Delta_{\mathcal{E:}}$}{Delta E})}
Drawing on Hyland's taxonomy of scientific hedges and boosters \cite{hyland1998hedging}, let $S_{\text{claim}}(c_k) \in [0, 1]$ and $S_{\text{evid}}(E^*) \in [0, 1]$ denote continuous epistemic commitment scores estimated by a token-level modality regressor trained on hedged (\emph{suggests, may, plausible}, $S \approx 0.2$), neutral empirical (\emph{observed, measured}, $S \approx 0.65$), and booster/causal (\emph{proves, definitively establishes, causes}, $S \approx 0.95$) constructions. Adjusting for the rhetorical section authority $w_{z(e)}$ of the supporting passages in $E^*$, we define the Epistemic Modality Delta as:
\begin{equation}
\Delta_{\mathcal{E}}(c_k, E^*) = S_{\text{claim}}(c_k) - \max_{e_r \in E^*} \Big( w_{z(e_r)} \cdot S_{\text{evid}}(e_r) \Big).
\label{eq:epistemic_delta}
\end{equation}
As illustrated in Fig.~\ref{fig:calibration}(a), when $\Delta_{\mathcal{E}}(c_k, E^*) > \tau_{\text{mod}}$ (where $\tau_{\text{mod}} > 0$ is the tolerable rhetorical variation band), the citing author commits \emph{epistemic overclaiming} (Class T1 or T2). Conversely, $\Delta_{\mathcal{E}} < -\tau_{\text{mod}}$ reflects conservative hedging, and $|\Delta_{\mathcal{E}}| \leq \tau_{\text{mod}}$ indicates calibrated epistemic fidelity.

\subsubsection{Structured Scope and Population Divergence (\texorpdfstring{$D_{\text{scope:}}$}{D scope})}
Let $\mathcal{F}(x) = (\mathbf{f}_{\text{pop}}(x), \mathbf{f}_{\text{inter}}(x), \mathbf{f}_{\text{metric}}(x), \mathbf{f}_{\text{cond}}(x))$ denote a four-slot frame representation (Population/Dataset, Intervention/Model, Outcome Metric, and Boundary Condition) extracted from text $x$, with slot set $\mathcal{S}_{\text{slot}} = \{\text{pop}, \text{inter}, \text{metric}, \text{cond}\}$. We define the scope divergence score $D_{\text{scope}}(c_k, E^*) \in [0, 1]$ as:
\begin{equation}
D_{\text{scope}}(c_k, E^*) = 1 - \prod_{s \in \mathcal{S}_{\text{slot}}} \left( \text{sim}_{\text{ont}}\!\left(\mathbf{f}_s(c_k), \mathbf{f}_s(E^*)\right) \right)^{\nu_s},
\label{eq:scope_div}
\end{equation}
where $\text{sim}_{\text{ont}}(\cdot, \cdot) \in [0, 1]$ measures asymmetric ontological subsumption (penalizing upward generalization from narrow source entities such as \emph{murine in vitro cells} to broad claim entities such as \emph{human patients}, while permitting valid downward specialization), and $\nu_s > 0$ ($\sum_s \nu_s = 1$) weights each slot.

\subsubsection{Neuro-Symbolic Quantitative Tolerance Check (\texorpdfstring{$\delta_{\text{num:}}$}{delta num})}
Let $\mathcal{Q}(c_k) = \{(v_q, u_q, \rho_q)\}_{q=1}^{Q}$ denote numerical tuples extracted from $c_k$, where $v_q \in \mathbb{R}$ is a standardized numerical value, $u_q$ is its canonical SI/statistical unit, and $\rho_q$ is its bound metric entity, and let $\hat{v}_q(E^*)$ be the aligned quantity extracted from prose or tables in $E^*$ \cite{herzig2020tapas,lu2023scitab}. The maximum relative numerical discrepancy is:
\begin{equation}
\delta_{\text{num}}(c_k, E^*) = \max_{q \in \{1, \dots, Q\}} \frac{|v_q - \hat{v}_q(E^*)|}{\max(|\hat{v}_q(E^*)|, \epsilon_0)},
\label{eq:num_check}
\end{equation}
where $\epsilon_0 = 10^{-8}$ prevents division by zero. A numerical distortion flag (Class T5) is triggered whenever $\delta_{\text{num}}(c_k, E^*) > \epsilon_{\text{tol}}$ (accounting for standard rounding precision) or when a statistical significance inequality (e.g., claiming $p < 0.05$ when $\hat{v}_p(E^*) = 0.054$) is violated.

\subsection{Split Conformal Risk Control and Severity Scoring}
To prevent uncalibrated neural posteriors from generating false accusations, we apply split conformal prediction \cite{sadinle2019least,angelopoulos2023conformal}. Let $\mathcal{D}_{\text{cal}} = \{(x_i, y_i^*)\}_{i=1}^{n_{\text{cal}}}$ be an exchangeable calibration set of expert-annotated claim--evidence pairs $x_i = (c_i, E_i^*)$ with ground-truth label $y_i^* \in \mathcal{Y}$. Using temperature-scaled posterior probabilities $\hat{P}_{\Theta, T}(y \mid x_i)$, we define the non-conformity score $s_i = s(x_i, y_i^*) = 1 - \hat{P}_{\Theta, T}(y_i^* \mid x_i)$. For a user-specified miscoverage risk level $\alpha \in (0, 1)$ (e.g., $\alpha = 0.05$), let $\hat{q}_{1-\alpha}$ be the finite-sample corrected empirical quantile over $\{s_1, \dots, s_{n_{\text{cal}}}\}$:
\begin{equation}
\hat{q}_{1-\alpha} = \text{Quantile}\!\left(\{s_1, \dots, s_{n_{\text{cal}}}\};\; \frac{\lceil (n_{\text{cal}} + 1)(1 - \alpha) \rceil}{n_{\text{cal}}}\right).
\label{eq:conformal_quantile}
\end{equation}
For a new test pair $x_{\text{test}} = (c_k, E^*)$, the split conformal prediction set is constructed as:
\begin{equation}
\mathcal{C}_{\alpha}(c_k, E^*) = \left\{ y \in \mathcal{Y} \;:\; 1 - \hat{P}_{\Theta, T}(y \mid c_k, E^*) \leq \hat{q}_{1-\alpha} \right\}.
\label{eq:conformal_set}
\end{equation}
By the exchangeability theorem of split conformal prediction \cite{sadinle2019least,angelopoulos2023conformal}, $\mathcal{C}_{\alpha}(c_k, E^*)$ satisfies the finite-sample marginal coverage guarantee $\mathbb{P}(y^*_{\text{test}} \in \mathcal{C}_{\alpha}(c_k, E^*)) \geq 1 - \alpha$ regardless of the underlying neural architecture.

Finally, for claims flagged as misaligned, we compute a bounded Misinterpretation Severity Score $\mathcal{S}_{\text{sev}}(c_k, E^*) \in [0, 1]$ to prioritize human editorial triage:
\begin{equation}
\begin{aligned}
\mathcal{S}_{\text{sev}}(c_k, E^*) = h_{\text{doc}} \cdot \sigma\Big( & \alpha_1 \hat{P}_{\Theta, T}(\textsc{Contradicted} \mid c_k, E^*) \\
& + \alpha_2 \max(0, \Delta_{\mathcal{E}}) + \alpha_3 D_{\text{scope}}(c_k, E^*) \\
& + \alpha_4 \min(1, \delta_{\text{num}} / \tau_{\text{num}}) \Big),
\end{aligned}
\label{eq:severity}
\end{equation}
where $\sigma(\cdot)$ is the logistic normalization function, $\alpha_1, \dots, \alpha_4 > 0$ are severity weights, and multiplying by the document access indicator $h_{\text{doc}} \in \{0, 1\}$ ensures that paywall-truncated sources ($h_{\text{doc}} = 0$) are never assigned high misinterpretation severity without full-text inspection.

\subsection{End-to-End Algorithmic Procedure}
Table~\ref{tab:algorithm_cedar} formalizes the complete end-to-end execution of the CEDAR framework for auditing a manuscript $D_{\text{cite}}$.

\begin{table}[t]
\centering
\footnotesize
\caption{Algorithm 1: End-to-End CEDAR Citation Verification Procedure}
\label{tab:algorithm_cedar}
\renewcommand{\arraystretch}{1.12}
\begin{tabularx}{\columnwidth}{Y}
\toprule
\textbf{Input:} Citing PDF $D_{\text{cite}}$, calibration quantile $\hat{q}_{1-\alpha}$, thresholds $(\tau_{\text{attr}}, \tau_{\text{mod}}, \tau_{\text{scope}}, \epsilon_{\text{tol}})$. \\
\textbf{Output:} Audit report $\mathcal{R} = \{\mathcal{V}(c_{i,k}, D_{\text{src}}^{(j)})\}$ for all inline citations. \\
\midrule
1: Parse $D_{\text{cite}}$ via GROBID \cite{lopez2009grobid} into sentence sequence $\mathcal{S}$, inline anchors $\mathcal{A}$, and bibliography $\mathcal{B}$. \\
2: Initialize audit report $\mathcal{R} \leftarrow \emptyset$. \\
3: \textbf{for each} inline citation anchor $a_i \in \mathcal{A}$ in focal sentence $s_{t(i)}$ \textbf{do} \\
4: \hspace{0.8em} Resolve $b_j \leftarrow \mu(a_i)$ via Crossref/OpenAlex/S2ORC \cite{lo2020s2orc}. \\
5: \hspace{0.8em} \textbf{if} $b_j$ does not resolve to an authentic catalog record \textbf{then} \\
6: \hspace{1.6em} Append $(\text{Anchor } a_i, \textsc{Phantom/Fabricated Citation})$ to $\mathcal{R}$; \textbf{continue}. \\
7: \hspace{0.8em} \textbf{end if} \\
8: \hspace{0.8em} Set horizon flag $h_{\text{doc}} \leftarrow 1$ if full text of $D_{\text{src}}^{(j)}$ is retrieved, else $h_{\text{doc}} \leftarrow 0$. \\
9: \hspace{0.8em} Index $D_{\text{src}}^{(j)}$ into $\mathcal{U}_{\text{prose}} \cup \mathcal{U}_{\text{table}} \cup \mathcal{U}_{\text{fig}}$ with section tags $z(e_r)$ (Eq.~\ref{eq:source_units}). \\
10: \hspace{0.5em} Compute optimal discourse window $C^*(a_i)$ via cohesion objective (Eqs.~\ref{eq:window_opt}--\ref{eq:cohesion_score}). \\
11: \hspace{0.5em} Decompose $C^*(a_i)$ into decontextualized atomic claims $\{c_1, \dots, c_m\}$ \cite{kamoi2023wice,min2023factscore}. \\
12: \hspace{0.5em} Filter claims attributed to $a_i$: $\mathcal{C}(a_i) \leftarrow \{c_k : A(c_k, a_i) \geq \tau_{\text{attr}}\}$ (Eqs.~\ref{eq:bipartite_logit}--\ref{eq:bipartite}). \\
13: \hspace{0.5em} \textbf{for each} attributed atomic claim $c_{i,k} \in \mathcal{C}(a_i)$ \textbf{do} \\
14: \hspace{1.3em} Retrieve top-$K_1$ candidates via SW-RRF score $S_{\text{SW-RRF}}(c_{i,k}, e_r)$ (Eqs.~\ref{eq:bm25}--\ref{eq:sw_rrf}). \\
15: \hspace{1.3em} Rerank top-$K_1$ via cross-encoder $g_{\text{cross}}$ to obtain top-$K_2$ rationale set $E^*$. \\
16: \hspace{1.3em} Compute calibrated NLI posteriors $\hat{P}_{\Theta, T}(y \mid c_{i,k}, E^*)$ over $y \in \mathcal{Y}$. \\
17: \hspace{1.3em} Compute $\Delta_{\mathcal{E}}(c_{i,k}, E^*)$ (Eq.~\ref{eq:epistemic_delta}), $D_{\text{scope}}(c_{i,k}, E^*)$ (Eq.~\ref{eq:scope_div}), and $\delta_{\text{num}}(c_{i,k}, E^*)$ (Eq.~\ref{eq:num_check}). \\
18: \hspace{1.3em} Construct conformal prediction set $\mathcal{C}_{\alpha}(c_{i,k}, E^*)$ using quantile $\hat{q}_{1-\alpha}$ (Eq.~\ref{eq:conformal_set}). \\
19: \hspace{1.3em} \textbf{if} $|\mathcal{C}_{\alpha}(c_{i,k}, E^*)| > 1$ \textbf{or} $h_{\text{doc}} == 0$ \textbf{then} \\
20: \hspace{2.1em} Set routing status $\leftarrow \textsc{Abstain / Human Expert Audit Required}$. \\
21: \hspace{1.3em} \textbf{else} \\
22: \hspace{2.1em} Assign deterministic verdict $\hat{y}_{i,k} \in \mathcal{C}_{\alpha}$ and taxonomy tag $\hat{\tau}_{\text{err}} \in \{\text{T1}, \dots, \text{T8}\}$. \\
23: \hspace{1.3em} \textbf{end if} \\
24: \hspace{1.3em} Compute severity $\mathcal{S}_{\text{sev}}(c_{i,k}, E^*)$ (Eq.~\ref{eq:severity}) and compile Verification Card into $\mathcal{R}$. \\
25: \hspace{0.5em} \textbf{end for} \\
26: \textbf{end for} \\
27: \textbf{return} $\mathcal{R}$. \\
\bottomrule
\end{tabularx}
\end{table}

\section{Critical Analysis of Large Language Models as Verification Engines}
\label{sec:llm_critique}

While instruction-tuned LLMs have demonstrated strong zero-shot capabilities in decomposing compound claims \cite{kamoi2023wice,min2023factscore} and generating natural-language explanations \cite{gao2023enabling,rashkin2023measuring}, deploying unconstrained generative LLMs as standalone citation verifiers introduces four documented failure modes:
\begin{enumerate}
\item \textbf{Parametric Memory Override (Prior Knowledge Intrusion):} Because LLMs encode broad scientific consensus within their pretrained parameters, an LLM prompted to verify whether a specific cited paper $D_{\text{src}}^{(j)}$ supports a generally accepted scientific statement frequently outputs \textsc{Supported} even when $D_{\text{src}}^{(j)}$ never investigated that proposition \cite{rashkin2023measuring,gao2023enabling}. Verification must strictly measure conditional attribution $P(c_k \mid E^*)$ rather than marginal plausibility $P_{\text{LLM}}(c_k)$.
\item \textbf{Long-Context ``Lost in the Middle'' Degradation:} Passing an entire 10,000-word manuscript directly into a long-context LLM prompt causes severe retrieval degradation. Empirical analyses by Liu et al.\ \cite{liu2024lost} prove that LLM recall drops sharply when relevant evidentiary sentences reside in the middle third of a long context window, justifying CEDAR's modular SW-RRF + cross-encoder retrieval hierarchy (Stages 7--8).
\item \textbf{Confabulatory Bridging and Sycophantic Framing Bias:} When presented with an authoritative-sounding citing sentence, generative models exhibit an affirmative bias toward validating the user-supplied claim and may confabulate unstated biochemical or mathematical steps to bridge a logical gap between $E^*$ and $c_k$.
\item \textbf{Fragility on Tabular and Numerical Inequalities:} Autoregressive token predictors frequently fail on multi-row table aggregations, confidence-interval overlap checks, and relative percentage comparisons \cite{lu2023scitab,akhtar2023reading}. CEDAR mitigates this vulnerability by delegating quantitative comparisons to symbolic inequality checks (Eq.~\ref{eq:num_check}) over structured table parses \cite{herzig2020tapas}.
\end{enumerate}

\section{Proposed Reproducible Evaluation Methodology}
\label{sec:evaluation}

\emph{Methodological Transparency Note:} In accordance with strict standards of scientific reproducibility, no synthetic experimental scores or fabricated benchmark tables are reported in this paper. Instead, we specify a rigorous, fully reproducible experimental design and evaluation protocol for benchmarking CEDAR against established baselines.

\begin{table*}[t]
\centering
\footnotesize
\caption{Proposed Multi-Dimensional Evaluation Protocol, Metrics, and Baseline Comparison Suite}
\label{tab:evaluation_metrics}
\renewcommand{\arraystretch}{1.14}
\begin{tabularx}{\textwidth}{L{2.7cm} L{3.3cm} L{4.2cm} Y}
\toprule
\textbf{Evaluation Dimension} & \textbf{Primary Quantitative Metrics} & \textbf{Mathematical Definition / Criterion} & \textbf{Target Baselines \& Statistical Testing Protocol} \\
\midrule
1. Atomic Claim Decomposition (RQ1) & Claim-Anchor Binding F1; Pronominal Decontextualization Rate & Harmonic mean of precision and recall over gold bipartite edges $(c_{i,k}, a_i)$ \cite{kamoi2023wice,min2023factscore}. & Single-sentence citance extraction ($w=0$); dependency-only clause splitting; paired bootstrap test ($B=10{,}000$, $p < 0.01$). \\
2. Multimodal Evidence Retrieval (RQ2) & Recall@$K$ ($K \in \{3, 5, 10\}$); Mean Reciprocal Rank (MRR) & $\text{MRR} = \frac{1}{|Q|}\sum_{q=1}^{|Q|} \frac{1}{\text{rank}_q}$ over gold prose, table, and caption rationales $E_{\text{gold}}^*$. & BM25 \cite{robertson2009probabilistic}; DPR \cite{karpukhin2020dense}; SciBERT \cite{beltagy2019scibert}; SPECTER \cite{cohan2020specter}; ColBERT \cite{khattab2020colbert}; unweighted RRF \cite{cormack2009reciprocal}. \\
3. Claim-Evidence Entailment \& Taxonomy (RQ3) & Macro-F1; SciFact Label+Rationale F1 \cite{wadden2020fact}; Taxonomy Cohen's $\kappa$ & Strict label + rationale overlap requiring $\hat{y}_{i,k} = y^*_{i,k}$ and $\text{IoU}(E^*, E_{\text{gold}}^*) \geq 0.5$. & scite intent classifier \cite{nicholson2021scite}; MultiVerS \cite{wadden2022multivers}; TaPas+NLI \cite{herzig2020tapas,lu2023scitab}; zero-shot RAG LLM \cite{gao2023enabling}; McNemar's test. \\
4. Epistemic Delta \& Numerical Audit (RQ3) & Overclaiming Detection F1; Mean Absolute Error on $\Delta_{\mathcal{E}}$; Numerical Recall & Precision/recall on T1/T2 modality shifts ($\Delta_{\mathcal{E}} > \tau_{\text{mod}}$) and T5 numerical errors ($\delta_{\text{num}} > \epsilon_{\text{tol}}$). & Standard 3-way NLI without modality regressor; LLM prompt-only numerical verification \cite{lu2023scitab,akhtar2023reading}. \\
5. Uncertainty \& Conformal Calibration (RQ4) & Empirical Coverage ($1-\hat{\alpha}$); Mean Set Size $|\mathcal{C}_{\alpha}|$; ECE; Brier Score & $\text{Cov} = \frac{1}{n_{\text{test}}}\sum_{i} \mathbb{I}(y_i^* \in \mathcal{C}_{\alpha}(x_i))$ vs.\ target $1-\alpha$, minimizing $|\mathcal{C}_{\alpha}|$ and false-accusation rate. & Uncalibrated softmax threshold; Temperature Scaling alone; Least Ambiguous Set-Valued Classifiers \cite{sadinle2019least,angelopoulos2023conformal}. \\
\bottomrule
\end{tabularx}
\end{table*}

\subsection{Benchmark Construction and Leakage-Free Dataset Splitting}
To evaluate all stages of CEDAR, we design a multi-tier evaluation benchmark combining four existing public datasets with a proposed expert-annotated full-text audit corpus:
\begin{enumerate}
\item \textbf{Existing Public Benchmarks:} (i) \emph{SciFact} \cite{wadden2020fact} and \emph{MultiVerS} full-text biomedical verification suites \cite{wadden2022multivers}; (ii) \emph{COVID-Fact} \cite{saakyan2021covidfact}; (iii) \emph{WiCE} sub-claim entailment dataset \cite{kamoi2023wice}; and (iv) \emph{SCITAB} \cite{lu2023scitab} and chart-verification datasets \cite{akhtar2023reading} for tabular and visual claim verification.
\item \textbf{Full-Text Citation Misinterpretation Audit Protocol:} Using open-access full-text pairs from S2ORC \cite{lo2020s2orc} across biomedicine, computer science, and materials science, domain experts annotate $(D_{\text{cite}}, a_i, D_{\text{src}}^{(j)})$ triples across four dimensions: gold atomic claims, gold multimodal rationale units $E_{\text{gold}}^*$, four-way entailment label $y^*$, and taxonomy class T1--T8. Inter-annotator agreement is quantified using Cohen's $\kappa$ and Krippendorff's $\alpha$, requiring adjudication by a senior third reviewer whenever annotators disagree.
\item \textbf{Citation-Graph Disjoint Splitting (Data Leakage Prevention):} A critical flaw in naive random splitting of citation datasets is \emph{source-document leakage}: if the same cited paper $D_{\text{src}}^{(j)}$ appears in both the training and test splits, neural encoders memorize source-specific facts. All splits (70\% train, 10\% validation, 10\% conformal calibration $\mathcal{D}_{\text{cal}}$, and 10\% test $\mathcal{D}_{\text{test}}$) must be partitioned by \emph{disjoint cited source documents and citation-graph connected components}.
\end{enumerate}

\subsection{Evaluation Metrics, Baselines, and Statistical Testing}
Table~\ref{tab:evaluation_metrics} summarizes the quantitative metrics and baseline comparisons across all five evaluation dimensions. In addition to standard retrieval metrics (Recall@$K$, MRR) and strict SciFact Label+Rationale Macro-F1 \cite{wadden2020fact}, we evaluate calibration via Expected Calibration Error (ECE), Brier Score, empirical conformal coverage $1-\hat{\alpha}$, and the \textbf{High-Confidence False Accusation Rate (HC-FAR)}: the fraction of genuinely supported citations erroneously assigned a singleton $\{\textsc{Contradicted}\}$ conformal prediction set. Statistical significance between CEDAR and baseline models is assessed using McNemar's test for classification decisions and paired bootstrap resampling ($B = 10{,}000$ iterations) with Holm--Bonferroni correction for multiple comparisons.

\section{Limitations, Threats to Validity, and Future Directions}
\label{sec:limitations}

\subsection{Technical Limitations and Threats to Validity}
\begin{enumerate}
\item \textbf{Commercial Paywalls and Epistemic Horizon Bias:} As established by Larivi{\`e}re et al.\ \cite{lariviere2015oligopoly}, a substantial share of global scholarly literature remains behind subscription paywalls. When $h_{\text{doc}} = 0$ (abstract-only access), CEDAR must abstain from issuing \textsc{Contradicted} or \textsc{Neutral} misinterpretation verdicts to avoid penalizing citations to methodology details or secondary tables located in the paywalled body.
\item \textbf{Subjectivity of Epistemic Hedging Boundaries:} As Hyland \cite{hyland1998hedging} and Jergas and Baethge \cite{jergas2015quotation} emphasize, scientific interpretation exists on a continuum. Distinguishing acceptable rhetorical synthesis from actionable overclaiming ($\Delta_{\mathcal{E}} > \tau_{\text{mod}}$) exhibits non-trivial expert disagreement, making conformal abstention and human-in-the-loop review indispensable.
\item \textbf{Complex Multi-Panel Visuals and Mathematical Proofs:} While TaPas \cite{herzig2020tapas}, SCITAB \cite{lu2023scitab}, and chart reasoning models \cite{akhtar2023reading} handle structured tables and standard plots, verifying claims grounded in multi-step symbolic theorem proofs, 3D protein structures, or electrophoretic Western blots remains an open computer-vision and formal-verification challenge.
\item \textbf{Absence of Executed Empirical Benchmarks in This Study:} Because this work presents a methodological framework and theoretical formulation without executing large-scale training runs over a newly curated full-text corpus, empirical validation of the exact recall and calibration gains of SW-RRF and $\Delta_{\mathcal{E}}$ remains an essential next step prior to production deployment.
\end{enumerate}

\subsection{Future Research Directions}
Promising extensions of the CEDAR framework include: (i) integrating formal theorem provers and computer algebra systems into Stage~10 to verify mathematical derivations; (ii) extending static document-pair verification into dynamic citation-network auditing that traces the genealogical origin of a distortion across multi-hop citation graphs \cite{greenberg2009citation,simkin2003read}; and (iii) embedding real-time, privacy-preserving verification assistants directly into LaTeX and word-processing environments to alert authors to modality drift and scope overgeneralization \emph{before} manuscript submission.

\section{Conclusion}
\label{sec:conclusion}

Verifying that a cited publication exists in a bibliographic database is fundamentally distinct from verifying that its empirical or theoretical content substantiates the claim for which it is cited. Misinterpreted but genuine citations bypass conventional editorial checks and citation-intent classifiers, allowing exaggerated, overgeneralized, or inverted assertions to accumulate unfounded authority across the scientific literature. In this paper, we critically audited existing citation verification paradigms and developed CEDAR, a mathematically grounded framework that unifies cohesion-penalized discourse scoping, bipartite atomic claim-to-anchor decomposition, section-weighted multimodal Reciprocal Rank Fusion over full text and tables, Epistemic Modality Delta ($\Delta_{\mathcal{E}}$) quantification, and split conformal risk control. By coupling fine-grained multimodal evidence retrieval with distribution-free uncertainty guarantees and transparent Verification Cards, CEDAR establishes a principled blueprint for human-in-the-loop scientific citation auditing.

\bibliographystyle{IEEEtran}
\bibliography{references}

\end{document}
